\documentclass{beamer}
% Možne velikosti pisav so: 8pt, 9pt, 10pt, 11pt, 12pt, 14pt, 17pt, 20pt

% Nastavitve za šumnike
\usepackage[T1]{fontenc}
\usepackage[utf8]{inputenc}
\usepackage[slovene]{babel}

% To potrebujemo, da bomo lahko delali zapiske za predavatelja
\usepackage{pgfpages}

% Treba je ločiti zapiske za predavatelja in prosojnice za občinstvo.
% To lahko naredimo v beamerju, glej
% https://gist.github.com/andrejbauer/ac361549ac2186be0cdb

% Primer tako narejenih prosojnic:
% http://math.andrej.com/wp-content/uploads/2016/06/hott-reals-cca2016.pdf

% Izberemo, ali bomo naredili samo prosojnice, samo zapiske ali oboje

\setbeameroption{hide notes}                        % samo prosojnice
%\setbeameroption{show only notes}                   % samo zapiski
%\setbeameroption{show notes on second screen=right}  % oboje

% Ali naj se \item-i pojavljajo zaporedoma, ali vsi naenkrat?
% Naj se pojavljajo zaporedoma:
% \beamerdefaultoverlayspecification{<+->}

% Minimalistični still
\mode<presentation>
\usetheme{default}
\usecolortheme{beaver}

% Malo manj dolgočasna pisava
\usepackage{palatino}
\usefonttheme{serif}

% Izklopimo navigacijske simbole, ker so neuporabni
\setbeamertemplate{navigation symbols}{}

% Aktiviramo oštevilčenje strani
\setbeamertemplate{footline}[frame number]{}

% Stil za zapiske
\setbeamertemplate{note page}{\pagecolor{yellow!5}\insertnote}

\usepackage{amsmath}
\usepackage{amssymb}
\begin{document}


% Naloga 1.3.1: Za dokument uporabite razred `beamer'.
% Ne dodajajte nastavitve za velikost pisave, kot je bila v datoteki `5-prosojnice.tex`.

% Naloga 1.3.2: vključite paket `predavanja'.

% Naloga 1.3.3: definirajte okolji `definicija' in `izrek'.
% Namig: z iskanjem po datotekah (Ctrl+Shift+F oz. Cmd+Shift+F) 
% poiščite niz `{definicija}' ali niz `{izrek}'.

{\theoremstyle{plain}
\newtheorem{izrek}{Izrek}
}


{\theoremstyle{definition}
\newtheorem{definicija}{Definicija}
}



% Naloga 1.3.4: pripravite naslovno stran z vsebino:
\title{Matematični izrazi in uporaba paketa \texttt{beamer}}
\subtitle{\emph{Matematičnih} nalog ni treba reševati!}
\institute{Fakulteta za matematiko in fiziko}
\date{}
\titlepage
% - datum: naj se ne izpiše; to dosežete z ukazom \date{}.
% Zgornje podatke nastavite z ukazi kot v dokumentih razreda `article`.
% Več o tem, kako se naredi naslovno stran, si preberite na naslovu na naslovu:
% https://www.overleaf.com/learn/latex/Beamer
% To stran preberite do vključno razdelka "Creating a table of contents".
% Ukaz `\titlepage` deluje podobno kot ukaz `\maketitle`, ki ste ga že srečali.

% Naloga 1.3.5: pripravite kazalo vsebine.
% 1. Naslov prosojnice, s kazalom vsebine naj bo "Kratek pregled"
% 2. S pomožnim parametrom `pausesections' (v oglatih oklepajih) 
%    določite, da naj se kazalo vsebine odkriva postopoma.
%    Poglejte, kako deluje ta ukaz.
% 3. Ker ni videti preveč lepo, pomožni parameter zakomentirajte.

\begin{frame}
  \frametitle{Kratek pregled}
  \tableofcontents[pausesections]
\end{frame}

\section{Paket \texttt{beamer}}
\input{1-paket-beamer.tex}

\section{Paketa \texttt{amsmath} in \texttt{amsfonts}}
\input{2-paketa-amsmath-amsfonts.tex}

\section[Matematika, 1. del\\\large{Analiza, logika, množice}]{Matematika, 1. del}

\section{Stolpci in slike}

\section{Paket \texttt{beamer} in tabele}

\section[Matematika, 2. del\\\large{Zaporedja, algebra, grupe}]{Matematika, 2. del}


\end{document}



